% Migrated from proposal/proposal.tex (5 October 2026) on 14 October 2026.
Health trackers process data that can reveal sensitive aspects of
a person's life. Research has examined their data practices and users' privacy
concerns, but the way providers describe possible downstream harms needs
closer study. This thesis proposes a content analysis of \highlight[orange!25]{the privacy
policies and the in-app interfaces of 50 widely used apps available in Germany; the interfaces
are inspected in cognitive walkthroughs.}
The analysis will distinguish descriptions of data flows, capabilities such as
inference, and consequences for users. It will examine how harms are presented
and how providers frame privacy protections, guided by literature-based harm
trees.
\highlight[orange!25]{Two large language models (LLMs) of different families will code all
text; a person will check every label on named harms, stance and protection framing without
seeing the models' labels, and a stratified sample of units, coded blind by the author, will
measure the models' errors and correct the reported figures.} The output will be an
evidence-linked account of harm communication and protection framing in the
captured material, not of data practices, legal compliance or user effects.\todo{Resolved
4 November 2026: the LLM method is specified in Sections~4.6.2 and~4.7.2 and summarised here;
store listings and websites are no longer coded surfaces (D17), so they left the first sentence.}

\noindent\textbf{Keywords:} privacy harms, health trackers, privacy communication, privacy-enhancing technologies, content analysis.
